\section{泛化：AI 的根本短板}

\subsection{模型泛化远不如人类}

Ilya 将泛化视为当前 AI 最根本的问题：

\begin{importantbox}{泛化的两个子问题}
\begin{enumerate}
    \item \textbf{Sample Efficiency}：为什么模型需要比人类多得多的数据才能学会同样的东西？
    \item \textbf{Teaching Difficulty}：为什么教模型某个东西比教人类要难得多？人类不需要 verifiable reward——一个 mentor 展示自己的思维方式，学生就能习得研究方法
\end{enumerate}
\end{importantbox}

\subsection{进化的角色}

对于视觉、听觉、运动等技能，进化可能提供了强大的先验（evolutionary prior）。人类手指灵活度远超机器人，5 岁小孩的汽车识别能力已足够自动驾驶。

但对于语言、数学、编程等\textbf{近代才出现的技能}，进化不太可能提供特定先验。然而人类在这些领域的学习效率仍然远超模型——这暗示人类可能拥有某种\textbf{更基本的、通用的学习算法}，而非仅靠领域特定的进化先验。

\begin{importantbox}{人类学习的鲁棒性}
人类的学习具有惊人的鲁棒性（robustness）：
\begin{itemize}
    \item 更少的样本
    \item 更少的监督（teenager 学开车不需要 verifiable reward）
    \item 极强的鲁棒性（在全新环境中仍能有效学习）
\end{itemize}
Ilya 认为存在某种尚未发现的\textbf{机器学习原理}可以实现类似的泛化能力——"I think it can be done. The fact that people are like that I think it's a proof that it can be done."
\end{importantbox}

\begin{knowledgebox}{人类神经元可能做了更多计算}
Ilya 提到一个可能的 blocker：人类神经元实际执行的计算量可能比我们认为的更多。如果这是真的，并且这对学习至关重要，那么实现人类级泛化可能需要更多算力。
\end{knowledgebox}

\subsection{本章小结}

泛化是当前 AI 最根本的短板。人类在进化不太可能提供先验的领域（如编程）仍展现出卓越的学习效率，暗示存在某种更基本的学习原理。Ilya 对此有想法，但由于竞争原因无法公开讨论。


